% Generated draft: requires mathematical and source review.
\documentclass[11pt]{article}
\usepackage[margin=25mm]{geometry}
\usepackage{fontspec}
\setmainfont{Latin Modern Roman}
\usepackage{amsmath,amssymb,mathtools,unicode-math}
\setmathfont{Latin Modern Math}
\usepackage{enumitem,needspace,xurl,hyperref}
\setlength{\parindent}{0pt}
\setlength{\parskip}{5pt}
\setlength{\emergencystretch}{3em}
\newcommand{\N}{\mathbb{N}}
\newcommand{\Z}{\mathbb{Z}}
\newcommand{\Q}{\mathbb{Q}}
\newcommand{\R}{\mathbb{R}}
\newcommand{\C}{\mathbb{C}}
\newcommand{\F}{\mathbb{F}}
\newcommand{\A}{\mathbb{A}}
\newcommand{\PP}{\mathbb{P}}
\newcommand{\E}{\mathbb{E}}
\newcommand{\eps}{\varepsilon}
\newcommand{\dd}{\,\mathrm{d}}
\begin{document}
\section*{UnsolvedMath ID: 2798}
\textbf{Source problem code:} catalogue entry
\subsection{Definitions and mathematical statement}
Let \(S\) be a closed oriented surface and \(\phi:S\to S\) the source's pseudo-Anosov diffeomorphism. Its mapping torus is
\[
M_\phi=(S\times[0,1])/((x,1)\sim(\phi(x),0)).
\]
The closed hyperbolic setting requires the negative-Euler-characteristic surface case; the closed-torus calculations mentioned separately in the source are not hyperbolic mapping-torus cases. Orientation and genus conventions should be read from the original formulation rather than silently extending the hyperbolic conclusion to torus bundles.

Write
\[
X_{\operatorname{SL}_2}(S)=\operatorname{Hom}(\pi_1(S),\operatorname{SL}_2(\C))/\!/\operatorname{SL}_2(\C).
\]
The smooth point is the character \([r]\) of a representation \(r\), not the homomorphism itself as a point of the quotient. The invariance condition is \([r\circ\phi_*]=[r]\); it does not require literal equality of homomorphisms. For each intended admissible integer \(n\), set \(q_n=\exp(2\pi i/n)\). The Kauffman bracket skein algebra \(\mathcal K^{q_n}(S)\) is the algebra of framed-link skeins in \(S\times[0,1]\), multiplied by stacking, with the parameter convention of the cited sources. That convention and the precise admissible orders of roots of unity are not explicitly reproduced in the supplied record.

Let \(\rho_n:\mathcal K^{q_n}(S)\to\operatorname{End}(V_n)\) be the intended finite-dimensional irreducible representation corresponding to \([r]\). If \(\Phi_n\) denotes the algebra action induced by \(\phi\), an explicit intertwiner convention is
\[
L_n\rho_n(a)=\rho_n(\Phi_n(a))L_n
\quad(a\in\mathcal K^{q_n}(S)),\qquad
L_n\in\operatorname{GL}(V_n),\quad\det L_n=1.
\]
This realizes the isomorphism on a common vector space so that trace and determinant are defined. The short source notation uses \(L\) for the resulting \(n\)-dependent family; it does not specify all representation-correspondence and transport choices. The background reports that the normalized absolute trace is independent of the remaining allowed scalar choice, under its representation hypotheses.

The exact requested asymptotic is
\[
\lim_{n\to\infty}\frac1n\log|\operatorname{Tr}L_n|
=\frac1{4\pi}\operatorname{vol}_{\mathrm{hyp}}(M_\phi),
\]
for every intended closed-surface pseudo-Anosov monodromy and invariant smooth character, not just a geometric-holonomy character. The source writes an ordinary limit over \(n\) with no parity or subsequence restriction. If a trace vanishes, its logarithm needs an explicit convention, such as the extended-real value \(-\infty\); no eventual nonvanishing theorem is assumed. The requested limit is not replaced by a limsup.

The supplied literature triage dated 2026-08-17 labels the row partially solved. It reports once-punctured-torus families of Bonahon--Wong--Yang, Pandey's four-punctured-sphere result under technical hypotheses, and figure-eight/Dehn-filling variants, while explicitly saying these do not settle the arbitrary closed hyperbolic case. Wang's closed-torus computation has zero limsup and is a zero-volume comparison, not the requested general formula. The background describes the conjecture as a toy version of Kashaev's volume conjecture, revisited by Baseilhac--Benedetti, and qualifies the skein/character correspondence as more or less established by the cited work. All qualifications, including puncture invariants in the separate punctured variants, remain source reports. The registry status is open; no general proof or independent current-status assessment is asserted here.
\subsection{Short English statement}
For a closed hyperbolic surface mapping torus, does the determinant-normalized skein intertwiner have trace growing exponentially at rate hyperbolic volume divided by \(4\pi\)? The source's punctured examples and zero-volume torus computations do not establish this general closed case. Root-order, representation and zero-trace conventions need exact alignment.
\subsection{Sources}
\begin{itemize}
\item Local UnsolvedMath registry, permanent ID 2798, source code catalogue entry, ``Kirby Problem 2.50,'' including supplied literature triage dated 2026-08-17.
\item Baykur, Kirby, and Ruberman, \emph{K3: A New Problem List in Low-Dimensional Topology}, Problem 2.50 (2026): \url{https://math.berkeley.edu/sites/default/files/surv-295-ruberman-watermarked-author-pdf.pdf}.
\item Francis Bonahon and Helen Wong, ``Representations of the Kauffman bracket skein algebra II: Punctured surfaces,'' \emph{Algebraic \& Geometric Topology} 17(6) (2017), 3399--3434, DOI 10.2140/agt.2017.17.3399, source label [BW17a].
\item Charles Frohman, Joanna Kania-Bartoszynska, and Thang L\^e, ``Unicity for representations of the Kauffman bracket skein algebra,'' \emph{Inventiones Mathematicae} 215(2) (2019), 609--650, DOI 10.1007/s00222-018-0833-x, source label [FKBL19].
\item Iordan Ganev, David Jordan, and Pavel Safronov, ``The quantum Frobenius for character varieties and multiplicative quiver varieties,'' \emph{Journal of the European Mathematical Society} 27(7) (2025), 3023--3084, DOI 10.4171/jems/1427, source label [GJS25].
\item Hiroaki Karuo and Julien Korinman, ``Azumaya loci of skein algebras,'' arXiv:2211.13700 (2022), source label [KK22].
\item Charles D. Frohman, Joanna Kania-Bartoszynska, and Thang T. Q. L\^e, ``Sliced skein algebras and geometric Kauffman bracket,'' \emph{Advances in Mathematics} 463 (2025), Paper 110118, 65 pp., DOI 10.1016/j.aim.2025.110118, source label [FKBL25].
\item Francis Bonahon, Helen Wong, and Tian Yang, ``Asymptotics of quantum invariants of surface diffeomorphisms I: conjecture and algebraic computations,'' arXiv:2112.12852 (2021), source label [BWY21]: \url{https://arxiv.org/abs/2112.12852}.
\item Rinat M. Kashaev, ``The hyperbolic volume of knots from the quantum dilogarithm,'' \emph{Letters in Mathematical Physics} 39(3) (1997), 269--275, source label [Kas97].
\item St\'ephane Baseilhac and Riccardo Benedetti, ``Quantum hyperbolic invariants of 3-manifolds with \(\operatorname{PSL}(2,\C)\)-characters,'' \emph{Topology} 43(6) (2004), 1373--1423, DOI 10.1016/j.top.2004.02.001, source label [BB04].
\item Francis Bonahon, Helen Wong, and Tian Yang, ``Asymptotics of quantum invariants of surface diffeomorphisms II: The figure-eight knot complement,'' arXiv:2203.05730 (2022), source label [BWY22].
\item Tushar Pandey, ``The Bonahon-Wong-Yang volume conjecture for the four-puncture sphere,'' arXiv:2311.13151, dated 2024 in the supplied bibliography, source label [Pan24].
\item Zhihao Wang, ``Kauffman bracket intertwiners and the volume conjecture,'' \emph{Algebraic \& Geometric Topology} 25(4) (2025), 2143--2177, source label [Wan25]: \url{https://doi.org/10.2140/agt.2025.25.2143}.
\item Tushar Pandey and Ka Ho Wong, ``Generalized Bonahon-Wong-Yang volume conjecture of quantum invariants of surface diffeomorphisms I: the figure eight knot complement,'' arXiv:2402.04483 (2024), source label [PW24].
\end{itemize}
\end{document}
